\section{The encoding of call-by-value  
$\lambda$-calculus  into the $\pi$-calculus }
\label{s:encoLpi}


\begin{figure}[t]
  \begin{mathpar}
    \begin{array}{rcl}
\encoI {V}  &\defeq& 
\bind p 
 \bout p y .    \encoIVa {V} y
\\[\myptD]
      \encoI {MN}  &\defeq&  
\bind p 
\res q \big(
\begin{array}[t]{l}
\encoIa M q
                         | %\\[\myptS]
\inp q y. \res r \big(
%\begin{array}[t]{l}
\encoIa N r | \\[\myptD]
\! \! \! \!\inp r w .\bout y {w',p'}.(\fwd {w'} w|\fwd {p'} p)\big)\big)
%\end{array} 
\end{array}\\[\mysp]
\encoIV {\lambda  x.  M}  &\defeq& 
\bind y 
 !\inp y {x,q}.\encoIa M q
\\[\myptD]
\encoIV x   & \defeq& 
\bind y 
 \fwd y x
   \end{array}
  \end{mathpar}
 \caption{The encoding of call-by-value  $\lambda$-calculus into \piI} % the internal \pc}
 \label{f:enc_internal}
%\Mybar
 \end{figure}

We recall here Milner's encoding of call-by-value $\lambda$-calculus, transplanted into
\piI. 

The core of any encoding of the  $\lambda$-calculus into a process calculus is the 
translation of function application. This  
becomes a particular form of parallel combination of two processes, the function, and its argument;  
 $\beta$-reduction is then modelled as a process interaction. 


The encoding of a $\lambda$-term is parametric over a name, that is, it is an abstraction.
This parameter may
be thought of as % the  \emph{location} of that term, or as 
its \emph{continuation}.  
A term that becomes a value signals so at its continuation name
and, in doing so, it grants access 
to the body of the value. 
Such body is replicated, so that the value may be
copied several times.
When the value is a function, its body can
receive two names: (the access to) its value-argument, and the
following continuation.
In the translation of application, first, the function is run, then
the argument; finally, the function is informed of its argument and 
continuation.
(Milner noticed \cite{Mil92s} that his
call-by-value encoding can be easily tuned so to mimic
forms of 
evaluation in which, in an application $MN$, the function $M$ is run
first, or the argument $N$ is run first, or function and argument are
run in parallel.  Here we follow  the first option,  the most common  one in
 call-by-value.)

As in OGS, so in \piI 
the encoding uses 
two kinds of names, 
the \emph{continuation names} $p,q,r,\dots$, and 
the \emph{variable names} $x,y,v,w\dots$.
%
Figure~\ref{f:enc_internal} presents the encoding.
The encoding function $\cbvSymb$, on values, is decomposed by means of the auxiliary
function  $\cbvSymb^\variant$; this will be useful later, e.g., in the encoding of
environments. 
\reviewComment{please type the different
components of the encoding; in particular say that V* is for values}
\reviewCommentDS{We have added an explanation}
Process   $\fwd ab$  represents a \emph{link} (sometimes called forwarder; 
for readability we have adopted the infix
notation $\fwd ab$ for the constant $\fwd{}{}$). It 
transforms all outputs at $a$ into outputs at 
$b$ (therefore $a,b$ are
names of the same sort); thus the body of $\fwd ab$ is replicated,
 unless $a$ and $b$ are {continuation names}:
%  (names such as
% % $p,q,r$ over which the encoding of a term is abstracted). 
% The
% definition of the constant $\fwd{}{}$ therefore is: 
 $$\begin{array}{rcl}
 \fwd{}{} &\defeq &
 \left\{ \begin{array}{ll}
 \abs{p,q} \inp p {x}.\bout {q} {y}
           . {\fwd{ y}{ x}}\\[\myptSmall] 
 \multicolumn{1}{l}{  
        ~\quad  \mbox{if $p,q$ are continuation names}
                    }              \\[\myptD]
 \abs{x,y} ! \inp x {z,p}.\bout y {w,q}
           .({\fwd{ q}{ p}}|\fwd w z)  
\\[\myptSmall]
 \multicolumn{1}{l}{  
           ~\quad \mbox{if $x,y$ are variable names}
 }\end{array} \right. 
 \end{array}
  $$
% $$
% \fwd{}{} \defeq 
% \abs{x,y} ! \inp x {\tilz}.\bout y {\tilw}.
% \fwd\tilw \tilz
% $$
% where, for $\tilz = z_1,\dots, z_n$ and 
%  $\tilw = w_1,\dots, w_n$, we write  $\fwd\tilw \tilz$ as an  abbreviation for 
% $\fwd {w_1}{z_1} | \ldots | 
% \fwd {w_n}{z_n}$.



The equivalence induced on call-by-value  $\lambda$-terms by their encoding into \piI 
coincides with 
Lassen's 
\emph{eager normal-form (enf) bisimilarity} \cite{lassentrees}.
%enf-bisimilarity. 

\begin{thmC}[\cite{DurierHS18}]
\label{t:adrien}
$\enca M \wbPI \enca N$ iff $M$ and $N$ are enf-bisimilar. 
\end{thmC} 

\DSOCTb

\noindent
In proofs about the behaviour of the \piI representation of $\lambda$-terms
we sometimes follow Durier et al.\ 
\cite{DurierHS18} \reviewComment{why the possessive?}
 \reviewCommentDS{It was an English typo, corrected}
and use an optimisation of Milner's encoding,
 reported in
Appendix~\ref{a:oc}.
%  together with %the results about 
% its correctness
% (Lemma~\ref{l:opt_sound}).  
The optimised encoding, 
indicated in the remainder as 
 $\qencV$,
   is 
obtained  from the  initial one 
$\cbvSymb$ 
by performing a few (deterministic) reductions,
at the price of 
 a more complex definition, in which a case analysis is made for the encoding of
 applications.
We only report here the  result about its behavioural correctness:

\begin{lemma}[correctness of the optimisation]
  \label{l:opt_sound} 
For all $M $, it holds that
 $\encoI M\exn\encoV M$.
Similarly, for all \AOGS configurations $\FF$, 
it holds that
 $
\encoConI \FF \exn\encoConV \FF$.
\end{lemma} 
\reviewComment{A bit confusing because you haven’t introduced this encoding}
\reviewCommentDS{well, i see no way of helping here, as we do not want to move the
  optimised encoding from the appendix to here. Hence i would not do anything here}

\noindent
The lemma is established by  algebraic
reasoning \cite{DurierHS18,Durier20}; its extension to encodings of configurations is straightforward.

\DSOCTe
